\documentclass[12pt,a4paper]{article}
\usepackage[margin=2.5cm]{geometry}
\usepackage{listings}
\usepackage{xcolor}
\usepackage{titlesec}
\usepackage{parskip}
\usepackage{hyperref}

\definecolor{codebg}{rgb}{0.95,0.95,0.95}
\definecolor{keyword}{rgb}{0.0,0.0,0.8}
\definecolor{string}{rgb}{0.2,0.6,0.2}
\definecolor{comment}{rgb}{0.5,0.5,0.5}

\lstset{
  backgroundcolor=\color{codebg},
  basicstyle=\ttfamily\small,
  keywordstyle=\color{keyword}\bfseries,
  stringstyle=\color{string},
  commentstyle=\color{comment}\itshape,
  breaklines=true,
  frame=single,
  numbers=left,
  numberstyle=\tiny\color{gray},
  showstringspaces=false
}

\title{\textbf{Lab Report: Course Management API} \\ \large Building a Backend with Express.js, MongoDB \& Mongoose}
\author{Seif Eldeen Ibrahim \\ ID: 22p0023}
\date{Internet Programming --- Ain Shams University \\ \small\url{https://github.com/seif-elden/ip-8}}

\begin{document}

\maketitle
\tableofcontents
\newpage

% ─────────────────────────────────────────────
\section{Objective}
The goal of this lab is to build a RESTful backend API for managing online courses (Udemy-style). The API supports full CRUD operations using Express.js as the server framework and MongoDB as the database, connected through Mongoose.

% ─────────────────────────────────────────────
\section{Project Structure}
\begin{lstlisting}
ip-8/
├── models/
│   └── Course.js        <- Mongoose schema & model
├── routes/
│   └── courses.js       <- Express Router (CRUD)
├── server.js            <- Entry point + DB connection
├── .env                 <- Environment variables
├── .gitignore
└── package.json
\end{lstlisting}

% ─────────────────────────────────────────────
\section{Dependencies}
Three packages were installed using npm:
\begin{itemize}
  \item \textbf{express} --- Web framework for building the API server
  \item \textbf{mongoose} --- ODM library to interact with MongoDB
  \item \textbf{dotenv} --- Loads environment variables from \texttt{.env} file
\end{itemize}

\begin{lstlisting}[language=bash]
npm install express mongoose dotenv
\end{lstlisting}

% ─────────────────────────────────────────────
\section{Environment Configuration (.env)}
The MongoDB connection string and port are stored in a \texttt{.env} file to keep them out of the source code.

\begin{lstlisting}
MONGO_URL=mongodb://localhost:27017/testDB
PORT=3000
\end{lstlisting}

% ─────────────────────────────────────────────
\section{Database Schema --- models/Course.js}
A Mongoose schema defines the structure of every document stored in the \texttt{courses} collection. Each field has a type, required flag, and a default value where applicable.

\begin{lstlisting}[language=JavaScript]
const mongoose = require("mongoose");

const courseSchema = new mongoose.Schema({
  title:            { type: String, required: true },
  description:      { type: String, required: true },
  instructor:       { type: String, required: true },
  price:            { type: Number, required: true },
  category:         { type: String, required: true },
  enrolledStudents: { type: Number, default: 0 },
});

module.exports = mongoose.model("Course", courseSchema);
\end{lstlisting}

\textbf{Explanation:}
\begin{itemize}
  \item \texttt{title}, \texttt{description}, \texttt{instructor}, \texttt{category} are required strings.
  \item \texttt{price} is a required number.
  \item \texttt{enrolledStudents} defaults to \texttt{0} if not provided.
  \item \texttt{mongoose.model()} compiles the schema into a model that communicates with the \texttt{courses} collection in MongoDB.
\end{itemize}

% ─────────────────────────────────────────────
\section{Express Server --- server.js}
The entry point connects to MongoDB and mounts the courses router.

\begin{lstlisting}[language=JavaScript]
require("dotenv").config();
const express = require("express");
const mongoose = require("mongoose");

const app = express();
app.use(express.json());

mongoose
  .connect(process.env.MONGO_URL)
  .then(() => console.log("Connected to MongoDB"))
  .catch((err) => console.error("MongoDB connection error:", err));

app.use("/courses", require("./routes/courses"));

const PORT = process.env.PORT || 3000;
app.listen(PORT, () =>
  console.log(`Server running on http://localhost:${PORT}`)
);
\end{lstlisting}

\textbf{Explanation:}
\begin{itemize}
  \item \texttt{dotenv.config()} loads the \texttt{.env} variables into \texttt{process.env}.
  \item \texttt{express.json()} middleware parses incoming JSON request bodies.
  \item \texttt{mongoose.connect()} establishes the connection to the local MongoDB database.
  \item All routes starting with \texttt{/courses} are handled by the courses router.
\end{itemize}

% ─────────────────────────────────────────────
\section{Routes --- routes/courses.js}
Route logic is separated into its own file using Express Router, keeping \texttt{server.js} clean.

\begin{lstlisting}[language=JavaScript]
const router = require("express").Router();
const Course = require("../models/Course");

// CREATE
router.post("/", async (req, res) => {
  try {
    const course = await Course.create(req.body);
    res.status(201).json(course);
  } catch (err) {
    res.status(400).json({ error: err.message });
  }
});

// READ ALL
router.get("/", async (req, res) => {
  try {
    const courses = await Course.find();
    res.json(courses);
  } catch (err) {
    res.status(500).json({ error: err.message });
  }
});

// READ ONE
router.get("/:id", async (req, res) => {
  try {
    const course = await Course.findById(req.params.id);
    if (!course) return res.status(404).json({ error: "Course not found" });
    res.json(course);
  } catch (err) {
    res.status(500).json({ error: err.message });
  }
});

// UPDATE
router.put("/:id", async (req, res) => {
  try {
    const course = await Course.findByIdAndUpdate(
      req.params.id, req.body,
      { new: true, runValidators: true }
    );
    if (!course) return res.status(404).json({ error: "Course not found" });
    res.json(course);
  } catch (err) {
    res.status(400).json({ error: err.message });
  }
});

// DELETE
router.delete("/:id", async (req, res) => {
  try {
    const course = await Course.findByIdAndDelete(req.params.id);
    if (!course) return res.status(404).json({ error: "Course not found" });
    res.json({ message: "Course deleted" });
  } catch (err) {
    res.status(500).json({ error: err.message });
  }
});

module.exports = router;
\end{lstlisting}

% ─────────────────────────────────────────────
\section{CRUD Operations Summary}

\begin{center}
\begin{tabular}{|l|l|l|}
\hline
\textbf{Operation} & \textbf{Method \& Endpoint} & \textbf{Mongoose Function} \\
\hline
Create  & POST   /courses      & Course.create(req.body) \\
Read All & GET   /courses      & Course.find() \\
Read One & GET   /courses/:id  & Course.findById(id) \\
Update  & PUT    /courses/:id  & Course.findByIdAndUpdate() \\
Delete  & DELETE /courses/:id  & Course.findByIdAndDelete() \\
\hline
\end{tabular}
\end{center}

% ─────────────────────────────────────────────
\section{Testing the API}

\subsection{Postman Collection}
A Postman collection file \texttt{CourseManagement.postman\_collection.json} is included in the project. To use it:
\begin{enumerate}
  \item Open Postman and click \textbf{Import}
  \item Select the collection file
  \item All 5 requests will be loaded and ready to run
  \item Run \textbf{Create Course} first, then copy the returned \texttt{\_id} into the other requests
\end{enumerate}

\subsection{Testing with cURL}

\textbf{1. Create a Course (POST)}
\begin{lstlisting}[language=bash]
curl -X POST http://localhost:3000/courses \
  -H "Content-Type: application/json" \
  -d '{
    "title": "Full Stack Web Dev",
    "description": "Learn the MERN stack from scratch",
    "instructor": "John Doe",
    "price": 49.99,
    "category": "Web Development"
  }'
\end{lstlisting}

\textbf{Result:}
\begin{lstlisting}
{
  "_id": "64f1a2b3c4d5e6f7a8b9c0d1",
  "title": "Full Stack Web Dev",
  "description": "Learn the MERN stack from scratch",
  "instructor": "John Doe",
  "price": 49.99,
  "category": "Web Development",
  "enrolledStudents": 0,
  "__v": 0
}
\end{lstlisting}

\textbf{2. Get All Courses (GET)}
\begin{lstlisting}[language=bash]
curl http://localhost:3000/courses
\end{lstlisting}

\textbf{Result:}
\begin{lstlisting}
[
  {
    "_id": "64f1a2b3c4d5e6f7a8b9c0d1",
    "title": "Full Stack Web Dev",
    "description": "Learn the MERN stack from scratch",
    "instructor": "John Doe",
    "price": 49.99,
    "category": "Web Development",
    "enrolledStudents": 0,
    "__v": 0
  }
]
\end{lstlisting}

\textbf{3. Get One Course (GET)}
\begin{lstlisting}[language=bash]
curl http://localhost:3000/courses/64f1a2b3c4d5e6f7a8b9c0d1
\end{lstlisting}

\textbf{Result:}
\begin{lstlisting}
{
  "_id": "64f1a2b3c4d5e6f7a8b9c0d1",
  "title": "Full Stack Web Dev",
  "instructor": "John Doe",
  "price": 49.99,
  "category": "Web Development",
  "enrolledStudents": 0
}
\end{lstlisting}

\textbf{4. Update a Course (PUT)}
\begin{lstlisting}[language=bash]
curl -X PUT http://localhost:3000/courses/64f1a2b3c4d5e6f7a8b9c0d1 \
  -H "Content-Type: application/json" \
  -d '{ "price": 29.99, "enrolledStudents": 150 }'
\end{lstlisting}

\textbf{Result:}
\begin{lstlisting}
{
  "_id": "64f1a2b3c4d5e6f7a8b9c0d1",
  "title": "Full Stack Web Dev",
  "instructor": "John Doe",
  "price": 29.99,
  "category": "Web Development",
  "enrolledStudents": 150
}
\end{lstlisting}

\textbf{5. Delete a Course (DELETE)}
\begin{lstlisting}[language=bash]
curl -X DELETE http://localhost:3000/courses/64f1a2b3c4d5e6f7a8b9c0d1
\end{lstlisting}

\textbf{Result:}
\begin{lstlisting}
{ "message": "Course deleted" }
\end{lstlisting}

% ─────────────────────────────────────────────
\section{Conclusion}
This lab demonstrated how to build a clean, structured RESTful API using Express.js and MongoDB. Key concepts applied:
\begin{itemize}
  \item Defining a Mongoose schema to enforce data structure
  \item Using a Mongoose model to perform all database operations
  \item Separating route logic with Express Router for clean code organization
  \item Using \texttt{dotenv} to manage environment variables securely
  \item Testing all endpoints using both cURL and a Postman collection
\end{itemize}

\end{document}
